\begin{abstract}
Logical anomaly detectors learn from synthetic defects, and every synthetic defect carries a label that says where the anomaly is. We show that this label decides where the detector's anomaly map points, and that changing it fixes a costly failure. On matched rearrangements of real MVTec LOCO pushpin trays, the state-of-the-art SALAD detector flags every faulty tray. Yet at a cutoff that finds 95\% of the faulty compartments, it also opens 11 of the 13 intact ones. Matched controls trace this to SALAD's convention of labelling an inserted part over every instance of its class. Retraining SALAD from scratch on identical inputs, changing only that label to the pasted footprint, removes the spread. On real pushpin defects, AU-sPRO rises from 30--44\% to 82--93\%. At an operating point fixed in advance, the detector finds 76--80\% of the defective compartments while opening fewer than 0.4 intact ones per tray, where the released detector's best policy finds 49\% and opens 1.1. Localization also improves in breakfast boxes. In screw bags, where extra parts are interchangeable, the class-wide convention remains the better choice. A synthetic label should point where the task locates the fault.
\end{abstract}
